% PFTF_alpha_소개.tex — 대학원생/입문자용 쉬운 설명서
% 규칙(논문생성규칙.md 명령 2): 한국 고등학생(인문계) 수준 수학만, 최대한 쉽게, 각주 풍부.
% 최종 PDF는 사용자 명시 요청으로 생성(xelatex + kotex).
\documentclass[11pt]{article}
\usepackage[margin=27mm]{geometry}
\usepackage{kotex}          % 한글 조판 (xelatex)
\usepackage{amssymb}        % ✓, ✗ 등 기호
\usepackage{setspace}
\usepackage{enumitem}
\usepackage[hidelinks]{hyperref}
\onehalfspacing
\setlength{\parskip}{4pt}

\title{\textbf{점을 이어 표면을 만들다}\\[4pt]
\large 3D 스캔에서 프린터까지, 그리고 ``더 똑똑한 자''는 정말 더 좋을까?\\[6pt]
\normalsize ── PFTF-alpha 프로젝트 쉬운 설명서 ──}
\author{\small (입문자·대학원생용 해설. 수식은 고등학교 수준까지만 씁니다.)}
\date{2026-07-25}

\begin{document}
\maketitle

\begin{abstract}
\noindent
3D 스캐너로 물건을 찍으면 컴퓨터 안에는 \emph{점들의 구름}만 남습니다. 이 점들을
이어 붙여 ``표면''을 만들어야 비로소 3D 프린터로 다시 출력할 수 있지요. 이 설명서는
(1) 점을 이어 표면을 만드는 고전적인 방법인 ``알파 복합체''를 소개하고,
(2) 우리 연구가 던진 질문 ---``방향까지 고려하는 더 똑똑한 자를 쓰면 표면을 더 잘
만들 수 있을까?''--- 를 설명하며,
(3) 그 답이 뜻밖에도 \emph{``아니오''}였다는 정직한 결과와, 그럼에도 건진 몇 가지를
이야기합니다. 어려운 수학은 쓰지 않고, 비유와 그림 같은 설명으로 풀어갑니다.
\end{abstract}

\section{우리가 풀려는 문제: 점 구름에서 표면으로}

여러분이 좋아하는 피규어를 3D 스캐너로 찍었다고 합시다. 스캐너가 컴퓨터에 남기는
것은 매끈한 조각상이 아니라, 표면 위 여기저기를 콕콕 찍은 \textbf{수만 개의 점}입니다.
이것을 \textbf{점군}\footnote{점군(點群, point cloud): 공간에 흩어진 점들의 모임.
스캐너·라이다·사진측량 등이 만들어 내는 가장 날것의 3D 데이터입니다.}이라고 부릅니다.

점군은 ``별자리''와 비슷합니다. 하늘에 흩어진 별(점)들을 보고 우리는 사자자리, 오리온자리
같은 ``모양''을 상상하지요. 컴퓨터도 똑같은 일을 해야 합니다. 흩어진 점들을 보고
``이 점과 저 점을 이으면 표면이 되겠구나''를 판단해서, 점들을 삼각형으로 촘촘히 이어 붙인
\textbf{메시}\footnote{메시(mesh): 수많은 삼각형을 이어 붙여 만든 표면. 3D 프린터·게임·영화
CG가 모두 이 삼각형 껍질로 물체를 표현합니다.}를 만들어야 합니다.

이 ``점 잇기''를 잘 해야 하는 이유는 분명합니다. 3D 프린터는 \textbf{빈틈없이 닫힌
껍질}\footnote{이런 성질을 ``워터타이트(watertight)'', 즉 ``물이 새지 않는''이라고
부릅니다. 껍질에 구멍이 뚫려 있으면 프린터는 안과 밖을 구분하지 못해 출력을 거부하거나
엉뚱하게 채웁니다.}을 요구하기 때문입니다.
\begin{itemize}[leftmargin=1.4em]
  \item 점을 \emph{너무 많이} 이으면 → 원래 떨어져 있어야 할 두 부분이 들러붙습니다(과충전).
  \item 점을 \emph{너무 적게} 이으면 → 표면에 구멍이 뚫려 프린트할 수 없습니다.
\end{itemize}
그러니까 ``적당히'' 이어야 하는데, 그 ``적당히''를 어떻게 정하느냐가 이 이야기의
핵심입니다.

\section{고전적인 해법: 알파 복합체와 ``공 굴리기''}

가장 널리 쓰이는 방법이 \textbf{알파 복합체}\footnote{알파 복합체(alpha complex):
1990년대에 정립된 고전적 표면 복원 도구. 뒤에 나오는 ``공''의 반지름이 그리스 문자
알파($\alpha$)라서 붙은 이름입니다.}입니다. 원리는 놀랄 만큼 직관적입니다.

반지름이 $\alpha$(알파)인 \textbf{공}을 하나 상상하세요. 이 공을 점군의 빈 공간
여기저기로 굴립니다. 공이 두 점 사이의 틈에 \emph{끼어 들어가지 못하면}, 그 두 점은
``가까운 사이''로 보고 선으로 잇습니다. 공이 \emph{쏙 들어가 버리면} 그 사이는 비워 둡니다.

즉 \textbf{공의 크기(알파)가 곧 ``얼마나 촘촘히 이을지''를 정하는 하나의 손잡이}입니다.
\begin{itemize}[leftmargin=1.4em]
  \item 공이 아주 크면 → 어지간한 틈은 다 막혀서, 점들이 뭉텅이로 들러붙습니다(과충전).
  \item 공이 아주 작으면 → 웬만한 틈에 다 들어가서, 표면이 숭숭 뚫립니다.
\end{itemize}
그래서 ``딱 맞는 공 크기''를 찾는 것이 표면 복원의 오랜 숙제였습니다.

\section{하나의 크기로는 안 되는 이유}

문제는, \textbf{좋은 공 크기가 장소마다 다르다}는 데 있습니다. 스캔한 물체는 보통
어떤 곳은 점이 빽빽하고(촘촘한 부분) 어떤 곳은 점이 듬성듬성합니다(성긴 부분).
\begin{itemize}[leftmargin=1.4em]
  \item 촘촘한 곳에 맞춰 공을 \emph{작게} 잡으면 → 성긴 곳이 구멍투성이가 됩니다.
  \item 성긴 곳에 맞춰 공을 \emph{크게} 잡으면 → 촘촘한 곳이 떡져 붙습니다.
\end{itemize}
하나의 공 크기로 두 마리 토끼를 다 잡을 수 없습니다. 그래서 연구자들은 ``장소에 따라
공 크기를 알아서 바꾸는'' 개선판들을 만들었습니다. 이 논문은 그중 잘 확립된 두 가지를
\textbf{비교 기준선}\footnote{기준선(baseline): 새 방법이 ``이것보다 나은가?''를 재기
위해 세워 두는 비교 대상. 육상으로 치면 넘어야 할 ``기록''입니다.}으로 삼습니다.

\begin{description}[leftmargin=0em,style=nextline]
  \item[\textmd{\textbf{B4 — 밀도 적응(density)}}]
  점이 촘촘한 곳에서는 공을 작게, 성긴 곳에서는 크게. 즉 \emph{그 자리의 점 간격}에
  맞춰 공 크기를 자동으로 조절합니다. ``장소에 맞춰 자의 눈금을 바꾸는'' 셈입니다.
  \item[\textmd{\textbf{B5 — 법선 이방성(normal-anisotropic)}}]
  여기에 한 가지를 더합니다. 표면은 어느 \emph{방향}으로 뻗어 있는지가 중요합니다.
  이 방법은 ``표면을 따라 옆으로는 잘 잇되, 표면을 \emph{가로질러} 위아래로는 함부로
  잇지 말라''고 방향에 따라 다르게 굽니다.\footnote{동그란 공 대신 표면에 납작하게 눌린
  ``럭비공(타원)''을 굴린다고 상상하면 됩니다. 옆으로는 길고 위아래로는 짧은 자를 쓰는
  것이지요. 표면이 향한 방향(법선)을 이용합니다. \emph{법선}이란 표면에 수직으로 세운
  화살표를 말합니다.}
\end{description}

\section{이 논문의 핵심 질문: ``더 똑똑한 자''는 더 좋을까?}

B5가 방향을 고려한다면, \textbf{방향을 더 정교하게} 고려하면 더 좋아지지 않을까요?
우리 연구실의 PFTF\footnote{PFTF: 우리 연구실이 여러 프로젝트에 공통으로 쓰는 이론
틀의 이름입니다. 핵심 아이디어는 ``점들 사이의 방향 관계''를 \emph{텐서}라는 수학
장치로 담는 것입니다. \emph{텐서}는 쉽게 말해 ``방향마다 크기가 다른 자'', 즉 방향
정보를 품은 표(表) 같은 것이라고 생각하면 됩니다.} 이론은 바로 그 ``더 똑똑한 자''를
제안합니다. 이 논문에서는 그것을 두 가지 형태로 시험합니다.

\begin{description}[leftmargin=0em,style=nextline]
  \item[\textmd{\textbf{P1 — 국소적으로 방향을 정교하게 굽는 자}}]
  각 점 주변에서 ``이웃 점들이 어느 방향으로 얼마나 촘촘한가''를 읽어, 그 자리에 딱 맞는
  \emph{찌그러진 자}\footnote{수학에서는 이것을 ``국소 SPD 계량''이라고 부릅니다. 겁먹을
  필요 없습니다. ``그 자리에서 방향마다 거리를 다르게 재는, 잘 찌그러진 자'' 정도로
  이해하면 충분합니다. SPD는 이 자가 ``뒤집히거나 터지지 않게'' 안전 조건을 만족한다는
  기술적 표시입니다.}를 만들어 냅니다. B5의 럭비공보다 더 자유롭게 기울고 늘어나는 자입니다.
  \item[\textmd{\textbf{P2 — 자신 없으면 안전하게 물러서는 자}}]
  P1이 ``이 자리는 내가 잘 모르겠다''고 느끼면, 무리하지 않고 검증된 B4로 돌아갑니다.
  똑똑함과 신중함을 섞은 셈입니다.
\end{description}

\noindent\textbf{그래서 질문은 한 문장입니다:}
\emph{이 정교하고 복잡한 ``방향 자''가, 이미 잘 작동하는 B4·B5보다 표면을 정말 더 잘
만들어 내는가?}

\section{정직하게 시험하기: ``동결 홀드아웃''}

좋은 질문에는 \textbf{공정한 시험}이 필요합니다. 여기서 이 논문이 특히 공을 들인
부분이 나옵니다.

시험이 공정하려면, \textbf{시험 문제를 미리 봐서는 안 됩니다.} 방법의 손잡이(공 크기 등)를
맞출 때 쓰는 ``연습 문제''와, 최종 성적을 매기는 ``시험 문제''를 철저히 갈라놓아야
합니다. 우리는 이것을 \textbf{동결 홀드아웃}\footnote{홀드아웃(hold-out): 미리 따로 떼어
숨겨 두는 ``시험 문제''. \emph{동결}(freeze)이란, 연습 단계에서 손잡이를 한번 정하면
시험 단계에서는 \emph{절대 손대지 않고 얼어붙게} 한다는 뜻입니다. 시험 보면서 답을
고치지 못하게 하는 것이지요.}이라고 부릅니다. 정답지(정밀한 기준 표면)는 오직 채점할
때만 쓰고, 방법이 표면을 고르는 순간에는 절대 보여 주지 않습니다.

게다가 채점도 깐깐하게 합니다. ``평균 점수가 조금 높다''는 것만으로는 합격시키지 않고,
\textbf{한 문제 한 문제마다} 기존 최강 기준선을 넘겼는지를 따집니다.\footnote{이를
``케이스별(casewise) 심사''라고 합니다. 반 평균이 1점 올랐다고 좋아하기 전에, ``쉬운
문제에서 벌고 어려운 문제에서 잃은 것''은 아닌지 문제별로 확인하는 것입니다. 훨씬
엄격한 기준입니다.}

\section{결과: 솔직하게 ``아니오''}

시험 결과는 뜻밖이었습니다.

\begin{quote}
\textbf{진짜 3D 모델들로 시험했더니, 정교한 ``방향 자''(P1·P2)는 기존 B5를 이기지
못했습니다.} 오히려 방향을 소박하게 다루는 B5가 가장 좋은 성적(표면 일치도
0.868)\footnote{여기서 점수는 ``복원한 표면이 정답 표면과 얼마나 겹치는가''를 0에서 1
사이로 나타낸 값(F-점수)입니다. 1에 가까울수록 좋습니다.}을 냈고, 우리의 어떤 정교한
방법도 그 벽을 넘지 못했습니다.
\end{quote}

왜 그럴까요? 실제 스캔한 매끈한 물체에서는 ``표면이 어느 방향인지''가 이미 꽤 뚜렷해서,
소박한 B5만으로도 충분했던 것입니다. 정교함을 더 얹는다고 없던 정보가 생기지는
않았습니다.\footnote{흥미롭게도, \emph{인공으로 만든 까다로운 예제}에서는 B5가 헤매고
우리 방법이 나아 보이기도 했습니다. 하지만 그것은 그 예제에만 있는 함정 때문이었고,
진짜 데이터에서는 그림이 뒤집혔습니다. ``인공 시험에서 좋아 보인다''가 ``현실에서
좋다''를 뜻하지 않는다는 교훈입니다.}

\section{그래도 건진 것들}

``아니오''가 답이라고 해서 빈손인 것은 아닙니다. 시험 과정에서 값진 것 세 가지를
건졌습니다.

\begin{description}[leftmargin=0em,style=nextline]
  \item[\textmd{\textbf{(1) 연결 방식을 아예 바꾸면 진짜로 나아진다 --- M1}}]
  지금까지의 방법들(B4·B5·P1·P2)은 \emph{어떤 점을 이을 수 있는지의 밑그림은 그대로
  둔 채}, 그 위에서 ``어느 선을 살릴까''만 골랐습니다. 그런데 밑그림에 이미 잘못 그어진
  선 --- 예컨대 \textbf{가까이 마주 본 두 면을 잘못 잇는 ``거짓 다리''}\footnote{거짓
  다리(false bridge): 원래 떨어져 있어야 할 두 표면이 가깝다는 이유로 잘못 이어져 버린
  연결. 얇은 책 두 권을 세워 두었는데 사이가 붙어 버려 한 권처럼 보이는 상황을 떠올리면
  됩니다.} --- 이 있으면, 점수만 다시 매겨서는 그 선을 지울 수 없습니다.
  M1\footnote{M1은 각 점에 \emph{가중치}를 주어 ``밑그림 자체''를 다시 그립니다.
  중요한 점은 크게, 덜 중요한 점은 작게 취급해 연결의 지도를 바꾸는 것입니다.}은
  아예 밑그림을 다시 그려서 이 문제를 정면으로 건드립니다. 그 결과 소박한 기준선 B4를
  \emph{확실히} 개선했습니다. ``점수 다시 매기기''가 아니라 ``지도 다시 그리기''가
  열쇠였던 것이지요.
  \item[\textmd{\textbf{(2) 확신이 없으면 ``안전하게 실패''한다 --- G4}}]
  자율주행차가 상황이 애매하면 스스로 속도를 줄이고 운전자에게 넘기듯,
  이 방법도 계산이 조금이라도 불확실하면 \textbf{``이건 안전하다''고 도장을 찍지 않고}
  보수적인 안전 경로로 물러섭니다.\footnote{이를 ``fail-closed(안전하게 실패)''라고
  합니다. 잘못된 결과를 조용히 ``괜찮다''고 통과시키는 일이 절대 없도록, 아홉 가지
  실패 상황을 일부러 시험해 확인했습니다.} 실수를 몰래 흘려보내지 않는다는 보증은,
  자동화된 프린트 공정에서 특히 값집니다.
  \item[\textmd{\textbf{(3) ``해도 안 되는 길''의 지도}}]
  거짓 다리를 \emph{지워서} 없애려는 온갖 시도가 하나같이 실패한다는 것도 꼼꼼히
  기록했습니다.\footnote{이유는 두 가지입니다. 첫째, 이미 그려진 밑그림에서 선을
  하나 지우면 그 자리에 새 구멍이 생겨 버립니다(다리를 없앴더니 구멍이 생김). 둘째,
  가장 어려운 예제에서는 두 면 사이 틈이 점 간격보다도 좁아서, \emph{그 자리만 들여다보는}
  어떤 방법으로도 둘을 구분할 정보 자체가 없습니다.} 뒤에 오는 연구자가 같은 막다른
  길에 시간을 낭비하지 않도록 하는, 그 자체로 쓸모 있는 ``실패 지도''입니다.
\end{description}

\section{왜 ``실패''를 논문으로 쓰는가}

성공만 논문이 되는 것은 아닙니다. 오히려 \textbf{정직한 음성(negative) 결과}는
과학에서 매우 중요합니다.
\begin{itemize}[leftmargin=1.4em]
  \item ``복잡하고 그럴듯한 방법이 늘 더 좋은 것은 아니다''를 \emph{엄밀한 시험으로}
  보였습니다. 소박한 방법(B5)이 이미 충분히 강했으니까요.
  \item ``인공 예제에서의 성공''에 속지 않으려면 \emph{진짜 데이터로, 문제를 미리 보지
  않고} 시험해야 한다는 방법론을 몸소 보였습니다.
  \item 남들이 같은 헛수고를 반복하지 않도록 ``여기까지 해 봤지만 안 되더라''를 지도로
  남겼습니다.
\end{itemize}
과장 없이 ``선행 기법을 이겼다''고 주장하지 \emph{않는} 대신, 무엇이 되고 무엇이 안 되는지를
정확히 말하는 것 --- 그것이 이 논문의 태도입니다.

\section*{한 문단 요약}
\addcontentsline{toc}{section}{한 문단 요약}
3D 스캔의 점 구름을 표면으로 바꾸는 ``알파 복합체''에서, 우리는 ``방향까지 정교하게
고려하는 자를 쓰면 더 좋아질까?''를 물었습니다. 진짜 데이터로 공정하게 시험한 답은
\textbf{``아니오''}였고 --- 소박하게 방향을 다루는 기존 방법(B5)이 이미 최고였습니다.
대신 \emph{연결의 밑그림 자체를 바꾸는 방법(M1)}은 기초 방법을 확실히 개선했고,
\emph{확신 없으면 안전하게 물러서는 안전장치(G4)}도 얻었으며, \emph{안 되는 길의 지도}도
남겼습니다. 정교함이 늘 이기는 것은 아니라는, 정직하고 유용한 결론입니다.

\end{document}
